\setcounter{chapter}{2}
\chapter{Architecture and Design}\label{chap:arch}
\graphicspath{{Chapter3/figures/}}

%==============================================================================
\pagestyle{fancy}
\fancyhf{}
\fancyhead[R]{\bfseries\rightmark}
\fancyfoot[R]{\thepage}
\renewcommand{\headrulewidth}{0.5pt}
\renewcommand{\footrulewidth}{0pt}
\renewcommand{\chaptermark}[1]{\markboth{{\chaptername~\thechapter.\ #1}}{}}
\renewcommand{\sectionmark}[1]{\markright{\thesection~ #1}}

% --- Self-healing figure helper --------------------------------------------
% Renders Chapter3/figures/<name>.{pdf,png} once it exists; until then shows a
% labelled placeholder naming the PlantUML source. Keeps the document
% compiling whether or not the diagram has been rendered yet.
%   #1 = base filename   #2 = width   #3 = caption   #4 = \label key
\makeatletter
\newcommand{\pumlfig}[4]{%
  \begin{figure}[H]\centering
    \IfFileExists{Chapter3/figures/#1.pdf}{\includegraphics[width=#2]{#1}}{%
    \IfFileExists{Chapter3/figures/#1.png}{\includegraphics[width=#2]{#1}}{%
      \fbox{\parbox{0.85\linewidth}{\centering\vspace{1.8em}%
        \textit{Diagram to be rendered from} \texttt{figures/#1.puml}.\\[0.4em]
        \footnotesize #3\vspace{1.8em}}}}}%
    \caption{#3}\label{#4}%
  \end{figure}}
\makeatother

\begin{spacing}{1.2}
%==============================================================================

\section*{Introduction}
This chapter details the architecture of CLMPilot: a structure designed
to satisfy the platform's requirements by construction. Instead of assembling the design feature by feature, we let a few
core forces drive the system's foundational choices.

3 primary forces drive the design. \textbf{Integration-first}: contract
documents remain entirely within the customer’s own storage, meaning the platform 
handles only their lifecycle management.
\textbf{Compliance by architecture}: the storage model natively provides the immutable,
reconstructable audit trails demanded by GoBD, GDPR, and ISO~27001.
\textbf{Multi-tenant isolation}: the architecture ensures that a single deployment
can securely serve multiple customer organisations, keeping each tenant's data
strictly isolated.

%--------------------------------------------------------------------
\section{Architectural drivers}
\label{sec:arch-drivers}

To solve the problem set out in Chapter~\ref{chap:analysis} and shape the
solution, we did start from the requirements where
any acceptable design must be satisfied, and then promoted to \emph{drivers}
that arbitrate every architectural decision in this chapter. The
non-functional requirements thus become 5 drivers. Each is a
constraint any acceptable design must meet, and each reappears in the later
chapters as the reason for a delivery-time decision.

\textbf{Auditability and immutability.} GoBD requires contract records to be
kept immutably and recoverable as of any past date for up to 10 years;
ISO~27001 adds change tracking (Section~\ref{sec:nfr-audit}); this is the
strongest of the 5 drivers, and it rules out any design in which history can be mutated,
or in which the audit trail is reconstructed after the fact; what happened
is the primary record, not a by-product of the current state.

\textbf{Integration-first.} A contract's document stays in cloud storage or any URL; 
the platform stores a reference, never the contents
(Section~\ref{sec:competitive-landscape}). So the architecture needs an
explicit integration boundary with interchangeable providers, and it must
keep working when a store is unreachable. This driver is fixed: it is
the product's differentiator.

\textbf{EU data residency.} The GDPR requires personal data to stay in the
European Union (Section~\ref{sec:nfr-compliance}). The stack hosts in the EU
by default, and every external dependency (email, the document stores, the
intelligence service) counts as a data-processing boundary.

\textbf{Multi-tenant isolation.} One deployment, many tenants
(Section~\ref{sec:actors}). Every command and query carries a tenant
context, and cross-tenant access is blocked by construction. Isolation that
depends on every query filtering by tenant will eventually be broken by one that does not.

\textbf{Maintainability under solo development.} One engineer builds and
runs the platform (Section~\ref{sec:methodology}). This driver selects
against complexity: a narrow stack, conventional patterns, automation over
manual work. Every added component adds to the load on a single engineer.

Availability and performance (Section~\ref{sec:nfr-perf}) matter too, but at
MVP scale they constrain operations more than structure.
Chapter~\ref{chap:delivered} takes them up where they can be measured.

%--------------------------------------------------------------------
\section{Why CQRS and Event Sourcing for CLM?}
\label{sec:why-cqrs-es}

The drivers of Section~\ref{sec:arch-drivers} now have to select the
platform's foundational pattern: how state is stored, changed and queried.
Because that choice is the hardest one to revisit later, we make it
explicitly rather than by habit: we fix the selection criteria, present
the candidate architectures, and justify the choice
(Section~\ref{sec:arch-choice}); the retained pattern is then presented in
detail (Section~\ref{sec:cqrs-es-pattern}).

\subsection{Architectural choice}
\label{sec:arch-choice}

4 criteria follow from the drivers. \emph{Compliance-grade history},
decisive: an immutable, attributable record from which any past state can be
rebuilt. \emph{Domain fit}: reads dominate, writes are rare and pass
multi-step validation. \emph{Tenant isolation by construction}: every record
and every access carries its tenant.

Table~\ref{tab:arch-dimensions} sets the established architectural style
against the one retained, on each of the 3 dimensions of the choice.

\begin{table}[H]
  \centering
  \caption{Established architectural style opposed on each dimension of the
           foundational choice.}
  \label{tab:arch-dimensions}
  \footnotesize
  \begin{tabularx}{\linewidth}{|>{\bfseries}l|X|X|}
    \hline
    Dimension & Established style & Style retained \\
    \hline\hline
    State persistence & Current-state persistence, the default of the classic
      3-tier application~\cite{Fowler2003}; history bolted on as a
      shadow-table audit log, a second write that can drift from the first. &
      Event~Sourcing: the append-only event store is the
      record~\cite{FowlerEventSourcing,Young2010}, current state a derivation,
      any past state a replay. \\
    \hline
    Domain modelling & Layered (N-tier) architecture with Active~Record / ORM
      mapping~\cite{Fowler2003}: one entity set serves queries and mutations
      alike. & CQRS: commands enforcing invariants and purpose-built read
      projections are separate pipelines~\cite{FowlerCQRS,Young2010}. \\
    \hline
    Communication & Request-response (REST/RPC) over a shared ACID database:
      synchronous, immediately consistent. & Event-driven: projections consume
      events asynchronously; eventual consistency on the read side. \\
    \hline
  \end{tabularx}
\end{table}

Two packaged architectures cut across the 3 dimensions: a
document-oriented store, current-state with a flexible schema, and a
workflow (BPM) engine, which models the approval process rather than the
contract.

The first criterion settles it. Current-state persistence cannot rebuild a
past state, and a shadow-table trail can be edited or diverge from the data
it describes (Section~\ref{sec:nfr-audit}). A document store inherits that
and has no transaction for a multi-step approval; a BPM engine adds a heavy
runtime and vendor lock-in, and still leaves the trail unsolved.
Event~Sourcing meets the criterion by construction, CQRS fits the read-heavy
workload, and both carry a tenant context on every event and every query.
The price is eventual consistency, long-lived event schemas and a learning
curve, accepted because state can be re-derived later, a missing history
cannot (Section~\ref{sec:cqrs-es-pattern}).

\subsection{CQRS and Event Sourcing}
\label{sec:cqrs-es-pattern}

A contract's life is a sequence of events: registered, document linked,
approval granted or refused, deadline passed, terminated. These are facts
that accumulate, not states that overwrite one another.
\textit{Event~Sourcing} stores that sequence and derives the current state
by replaying it, instead of mutating state in
place~\cite{FowlerEventSourcing,Young2010}. \textit{CQRS} is its
companion~\cite{FowlerCQRS,Young2010}: it splits the write model (commands,
which validate rules and emit events) from the read model (queries, which
read purpose-built projections). The two are independent but reinforce each
other~\cite{MicrosoftCQRS,MicrosoftEventSourcing}.

This pairing satisfies the auditability driver. Under
Event~Sourcing the immutable, time-stamped, actor-attributed record that
GoBD and ISO~27001 demand (Section~\ref{sec:nfr-audit}) \emph{is} the event
log: there is no separate audit table to maintain or to fall out of sync.
\textit{Who approved this contract, when, and against which version} is
answered by replaying its events.

CQRS also fits the workload. Reads dominate and fetch whole aggregates: the
dashboard counts expiring contracts and sums value by quarter. Writes are
rarer and guarded: an approval or an edit must pass validation. Splitting
them lets each read model be its own projection~\cite{Fowler2003}, so a
heavy query never blocks a write, and it expresses the domain in the terms
of Domain-Driven Design~\cite{Evans2003,Vernon2013}.
Figure~\ref{fig:cqrs-es} shows both paths.

\pumlfig{cqrs-es}{0.90\textwidth}{The CQRS and Event Sourcing flow}{fig:cqrs-es}

Figure~\ref{fig:cqrs-seq} shows the same flow as a sequence, for
registering a contract. The command path is short and synchronous:
validate, append an event, answer the caller. The reactions come after: a
projection updates the read models, and a reactor turns the new event into a
follow-up command that creates the contract's deadlines. That is the reactor
pattern, and the point where eventual consistency becomes visible.

\pumlfig{cqrs-sequence}{\textwidth}{Registering a contract as a
sequence}{fig:cqrs-seq}

The costs anticipated in Section~\ref{sec:arch-choice} are now concrete.
Read models are eventually consistent: a new contract
may show in its detail view a moment before it shows in a dashboard count,
which for CLM is acceptable. Events are a long-lived schema, so they need a
versioning convention. And the pattern has a learning curve, mitigated
here because an existing framework can supply the mechanics; the choice of
that framework is the subject of the next section.


%--------------------------------------------------------------------
\section{The comby framework}
\label{sec:comby}

The framework that carries the patterns of
Section~\ref{sec:why-cqrs-es} did not appear for this project: it comes
out of Gradient~Zero's own history. The company builds compliance-first
products for regulated sectors
(Section~\ref{sec:areas-expertise}), and each new product kept running
into the same problems: an audit trail to guarantee, approval workflows
to model, temporal queries to answer: infrastructure that would
otherwise be reimplemented, or imported from an open-source project, once
per product. Its response was to consolidate that recurring
infrastructure into an in-house Go framework, maintained as a strategic
asset and reused across the portfolio.

comby is that framework. It supplies the CQRS/ES building blocks
(command and query buses, an append-only event store, a projection runner,
a reactor scaffold, and a message broker) and a set of \textit{defaults}:
tenants, accounts and identities, group-based access control, invitations,
asset storage, webhooks, and an auto-generated REST API with an OpenAPI
specification. CLMPilot runs on comby~v3.

\subsection{Role in the architecture}
\label{sec:comby-role}

CLMPilot is a thin layer of domain meaning over a thick layer of framework
mechanics (Figure~\ref{fig:comby-layers}). comby owns the \emph{how}:
dispatching and authorising a command, persisting and publishing an event,
rebuilding a projection, isolating a tenant. CLMPilot owns the \emph{what}:
what a contract is, when an approval advances, how a deadline raises
reminders. That split is the maintainability driver in practice.

\pumlfig{comby-layers}{0.80\textwidth}{Layering of CLMPilot on comby.}{fig:comby-layers}

What CLMPilot inherits is large: tenancy, accounts, identities, groups,
invitations, asset storage, and the whole authentication layer. The REST
API and its OpenAPI document are generated from the registered commands and
queries, so the frontend uses a generated, type-safe client. What CLMPilot
adds is small: the 4 core domains of Section~\ref{sec:domain-model}, the
handful of enterprise contexts that extend them on the same pattern
(Section~\ref{sec:extension-contexts}), and 2 cross-cutting pieces: a
middleware that blocks archived tenants, and a boot-time reactor that seeds
the system administrators.
Figure~\ref{fig:backend-packages} shows the same split as packages: the
domains register on comby, the satellites depend on Contract, and the
integration and internal packages keep the external services and
cross-cutting concerns apart.

\pumlfig{backend-packages}{\textwidth}{Packages structure of the backend}{fig:backend-packages}

\subsection{Justification of choice}
\label{sec:comby-choice}

Choosing comby is a build-versus-buy decision, settled by fit rather than
feature count. comby is a strategic asset of Gradient~Zero, not a
convenience library (Section~\ref{sec:gradient-zero}): CLMPilot exists
partly to exercise it, so its authors are on hand to support it.
Table~\ref{tab:framework-comparison} sets it against the alternatives that
were considered.

\begin{table}[H]
  \centering
  \caption{CQRS and Event~Sourcing platforms considered for CLMPilot,
           assessed against the project's decisive criteria.}
  \footnotesize
  \begin{tabularx}{\linewidth}{|>{\bfseries}l|l|X|X|}
    \hline
    Option & Ecosystem & CQRS/ES core \& built-in tenancy, identity, RBAC
      & Support and strategic fit \\
    \hline\hline
    comby~v3 & Go & Command/query buses, append-only event store,
      projection runner and reactors, \emph{plus} tenant, account,
      identity, group and invitation defaults out of the box. & Authored
      in-house (Section~\ref{sec:areas-expertise}); the framework team
      supports the product directly; reused across the company's
      portfolio. \\
    \hline
    EventStoreDB + Go~\cite{EventStoreDB} & Go & A mature event store
      only; the command and
      query buses, projection wiring, multi-tenancy, identity and RBAC
      must all be hand-built. & External product; the entire application
      framework remains the project's burden. \\
    \hline
    Axon Framework~\cite{AxonFramework} & Java/JVM & Mature CQRS and
      Event~Sourcing, but no
      multi-tenant identity layer; would require a JVM service. &
      Open-source; adopting it would split the stack across 2 languages.
      \\
    \hline
    EventFlow~\cite{EventFlow} & .NET & CQRS/ES primitives for the .NET
      ecosystem; tenancy
      and identity not provided. & Open-source, but the wrong ecosystem
      for a Go house. \\
    \hline
    Bespoke (in Go) & Go & Full control over every primitive, at the cost
      of re-implementing all of the above. & No external dependency, but
      unmaintainable by one engineer alongside the product itself. \\
    \hline
  \end{tabularx}
  \label{tab:framework-comparison}
\end{table}

3 points decide it: comby is Go-native, so the backend stays one
language. Its CQRS/ES primitives are proven, so they need not be rebuilt.
And its defaults supply the biggest saving: the tenancy, identity, RBAC
and invitation logic that every other option would leave to the project.
EventStoreDB and a bespoke build leave the whole framework to build; Axon
and EventFlow are the wrong ecosystem.

%--------------------------------------------------------------------
\section{Domain model overview}
\label{sec:domain-model}

The requirements of Chapter~\ref{chap:analysis} decompose into 4 core
\textit{bounded contexts}, each with its own aggregate, events
and rules. We give them here at the conceptual level (responsibilities and
relationships) and their realisation in Chapter~\ref{chap:realisation}.
The enterprise capabilities of Section~\ref{sec:fr-enterprise} add a handful
of further contexts that follow the same pattern; they are introduced at the
end of this section so the core model stays readable.

\subsection{Bounded contexts}
\label{sec:bounded-contexts}

Figure~\ref{fig:domains} shows the 4 bounded contexts and how they
relate. Contract is central; Approval,
Deadline and DocumentLink each reference it. All 4 sit inside a tenant
boundary: every instance belongs to one tenant, and no relationship
crosses it. Identity and the audit log are absent here: they are
platform concerns owned by the framework (Section~\ref{sec:comby-role}), not
business contexts.

\pumlfig{domain-model}{0.85\textwidth}{The 4 bounded contexts, within the tenant boundary}{fig:domains}

\subsection{Aggregate responsibilities}
\label{sec:aggregate-responsibilities}

\textbf{Contract} is the record of a contract's metadata and status:
counterparty, type, value, dates, renewal terms. It does not hold the
document; that is a DocumentLink. Every other context points at it, and it
anchors the audit trail.

\textbf{Approval} is the multi-step review a contract passes before it goes
active. It is its own aggregate, not a status field on Contract. That gives
the approval chain its own lifecycle and its own immutable record of who
decided what, and when.

\textbf{Deadline} tracks cancellation windows, renewal and expiry dates, and
the reminders before them. It is a first-class aggregate, so one
contract can carry several independent deadlines. That is how the platform
answers the missed-deadline problem: each deadline is an object the contract
owns, not just a field on it.

\textbf{DocumentLink} is the integration boundary turned into a domain
object. It records that a contract is embodied by a file in an external
store (Drive, SharePoint, or any URL) without ingesting the file. One
aggregate serves all providers, through provider-specific adapters.

\subsection{Behavioural views: lifecycle and workflow}
\label{sec:behavioural-views}

2 contexts are clearest through their behaviour.
Figure~\ref{fig:contract-states} is the Contract state machine. A contract
starts as a \textit{draft}, goes \textit{active} once approved, may turn
\textit{expiring} near its end date and then \textit{expire}; an active
contract can instead be \textit{terminated} early. Both paths end in
\textit{archived}. Each transition is a command or a time-driven event, so the diagram
is a compact summary of a contract's audit trail.
\\Soft delete (tombstone) and restore are available from any state and are omitted
 here for clarity.

\pumlfig{contract-lifecycle}{0.5\textwidth}{State machine of the Contract
aggregate}{fig:contract-states}

Figure~\ref{fig:approval-activity} shows the approval workflow as activity.
The Legal~Manager starts it against a template of ordered steps. Each
approver either advances it or interrupts it by rejecting or requesting
changes; only when every step is approved does the contract go active. The
diagram makes the key rule explicit: a single rejection returns the whole
contract to its initiator.

\pumlfig{approval-activity}{0.4\textwidth}{Activity diagram of the approval
workflow}{fig:approval-activity}

\subsection{Extension contexts}
\label{sec:extension-contexts}

The enterprise capabilities of Section~\ref{sec:fr-enterprise} did not change
the core model; they added new contexts that follow the same aggregate-and-events
discipline, referencing \texttt{Contract} exactly as the core satellites do.

\textbf{Signature} manages the electronic-signature lifecycle of a contract, tracking 
the envelope sent to external counterparties and individual signer outcomes. Modelled
similarly to \emph{Approval}, this is built as a distinct multi-step aggregate rather
than a simple flag on \texttt{Contract}. \textbf{Obligation} handles post-signature 
commitments and their recurrence, to track what the organisation is required to do once a
contract goes active. \textbf{CustomField} enables per-tenant schema extensions, 
allowing organisations to record their own custom attributes against contracts without
altering the core aggregate. \textbf{Reporting and analytics}, finally, are handled by a family of read projections driven 
by the events that other contexts already emit. 

These capabilities arrived as new contexts and projections, not as edits to the 4
core aggregates: the decomposition let the enterprise features extend the model at its seams,
without disturbing its core.

%--------------------------------------------------------------------
\section{System topology}
\label{sec:topology}

The structure has 2 views: a \textit{component} view (which services
exist, and through which interfaces) and a \textit{deployment} view (where
the code runs and the data is stored). Both describe the design; the operational
detail (the host, the reverse proxy, the backups, the release process)
is Chapter~\ref{chap:realisation}.

\subsection{Component view}
\label{sec:component-view}

The services connect through a small set of named interfaces
(Figure~\ref{fig:component-view}). The
backend provides one interface, the generated REST/OpenAPI contract, which
the frontend consumes through its generated client. The backend requires
3: \texttt{IntelligenceClient} from the intelligence service,
\texttt{Email} from the mail provider, and \texttt{DocumentStore} behind
which the cloud storage adapters sit.

\pumlfig{component-view}{0.5\textwidth}{CLMPilot component view}{fig:component-view}

\subsection{Deployment topology}
\label{sec:deployment-topology}

The deployment runs as a single EU-hosted, containerised stack
(Figure~\ref{fig:topology}) of 3 runtime services: the \textbf{backend},
the \textbf{frontend} (a single-page application running in the user's
browser) and the \textbf{intelligence} service (backed by a large
language model). We name the parts by their role; the concrete products
behind them are a realisation decision
(Chapter~\ref{chap:realisation}). The services are backed by 4 stores.
A \textbf{relational database} is the system of
record (the event store and the command store) and the source every read
model is rebuilt from. A \textbf{cache database} holds sessions and
accelerates read models. A \textbf{message broker} carries events to
projections and reactors and is designed to distribute them across
backend instances when the platform scales
horizontally. An \textbf{object store} keeps exports, signed
documents and other assets. One reverse proxy terminates TLS in front of
them.

\pumlfig{topology}{0.9\textwidth}{CLMPilot deployment topology}{fig:topology}

Network isolation is built into the topology: only the reverse proxy is reachable
from outside, while the intelligence service and the stores are reachable only from
the backend, so the most sensitive parts are not just access-controlled but
unaddressable. Outward, the system depends on an email provider, on the
users' own cloud storage for documents, and on an e-signature provider,
each behind a named interface so the concrete provider is interchangeable;
ERP integration is anticipated in the design but not yet built.

%--------------------------------------------------------------------
\section{Security and compliance model}
\label{sec:security}

Security is a property of the architecture, not a subsystem beside it. The
comby framework supplies the mechanisms (identity, access control and tenant
isolation) and CLMPilot inherits it (Section~\ref{sec:comby-role}). We
describe the identity and access model it provides, the layered pipeline that
enforces that model on every request, the defence in depth that keeps tenants
apart, and finally compliance, which largely restates
structure already in place.

\subsection{Identity and access model}
\label{sec:identity-model}

The model has 4 levels from the comby defaults, shown in
Figure~\ref{fig:identity-model}. A \textbf{tenant} is a customer
organisation. An \textbf{account} is a person's login, which may hold an
\textbf{identity} in more than one tenant: the same person can act in
several organisations without sharing a password. A \textbf{group} is a role
within a tenant (Legal, Finance, Admin); a \textbf{permission} is the right
to run a given command or query. Access control is role-based and always
scoped to a tenant.

\pumlfig{identity-model}{\textwidth}{The identity and access model}{fig:identity-model}

Above the per-tenant model sits the platform \textbf{system administrator}:
a small, seeded set of operator accounts in a reserved system tenant, able
to provision, inspect and archive tenants through a cross-tenant interface.

\subsection{Authentication and the request security pipeline}
\label{sec:security-pipeline}

Authentication establishes who is calling, and does so once, at the edge. A
user signs in with an email and a password (held only as a salted hash) 
and a second factor, a one-time code sent by email. What the browser then carries is a
cookie backed by an opaque, server-side session, not a self-contained bearer
token: the session holds no claims a client could read or replay, and it can
be revoked at will. comby's account and session handling covers this.

From that point, authorisation is a stage every request passes through: the
framework enforces it, so the application never has to invoke it explicitly. The framework resolves the
caller's identity and tenant once, from the session, then checks each command
or query against the caller's permissions \emph{before} it reaches a domain
handler. Figure~\ref{fig:security-pipeline} shows the 2 halves: identity is
established at the edge, and authorisation runs at dispatch, within the
caller's tenant. A request beyond the caller's authority, or one aimed at
another tenant, is refused at the boundary; and because the tenant context
is taken from the session rather than from the request, the ordinary API
has no way to address another tenant's data at all.

\pumlfig{security-pipeline}{0.88\textwidth}{The request security pipeline}
{fig:security-pipeline}

The authorisation check itself is a short, ordered cascade in which the first
rule that resolves decides the outcome (Figure~\ref{fig:authz-decision}).
Internal operations (a seed at boot, a reactor acting on an event) run
without a caller and are allowed; a command may carry a custom rule for a case
that role membership alone cannot express; a platform administrator is allowed
across tenants; any request whose target tenant differs from the caller's is
denied outright; and otherwise access is granted only if one of the caller's
tenant groups carries the matching permission. The framework offers a finer,
workspace-scoped grant beneath the tenant as well, which CLMPilot does not yet
exercise.

\pumlfig{authz-decision}{0.95\textwidth}{Authorisation process of commands and queries}
{fig:authz-decision}

\subsection{Defence in depth and tenant isolation}
\label{sec:defense-in-depth}

No single control is trusted to be enough: the same request crosses several
independent layers, each of which would hold if an outer one were bypassed
(Figure~\ref{fig:security-layers}). The network perimeter exposes only the
reverse proxy and keeps the stores on an internal-only network, unaddressable
from outside (Section~\ref{sec:deployment-topology}). Authentication and
per-tenant authorisation are the 2 layers just described. Beneath those 2
layers is the tenant isolation that all of the compliance guarantees depend on,
and at the core the immutable event log, which serves as both the system of
record and the audit trail.

\pumlfig{security-layers}{0.8\textwidth}{Defence in depth}{fig:security-layers}

Tenant isolation today is enforced in the application: every command, query
and event carries a tenant context, and a single logical event store is
partitioned by that context rather than split into physical per-tenant streams
(Section~\ref{sec:identity-model}). That isolation is therefore only as strong
as the authorisation layer above it, which refuses cross-tenant requests.
comby~v3 adds a deeper reinforcement, which CLMPilot does not yet activate: row-level security in the database
itself, so that the data layer would refuse a cross-tenant read even if the
application layer were wrong.

\subsection{Compliance positioning}
\label{sec:compliance-positioning}

The compliance controls follow directly from the architecture. GoBD's
immutable-retention requirement is met by the event store
(Section~\ref{sec:why-cqrs-es}): the audit trail \emph{is} the log. The
GDPR's residency requirement is met by EU hosting; its data-minimisation
principle, by integration-first design, since the platform never holds the
documents at all. ISO~27001's access-control and change-tracking
expectations are met by per-tenant RBAC and, again, the log.

One tension remains: an append-only store versus the right to
deletion. The platform reconciles them by \textit{tombstoning}: a deletion
is itself an event that strips the personal data but keeps the structure and
ordering of the history. The record that something happened is preserved,
while the personal data it contained is removed. That is only possible because
corrections here are events, not in-place edits.

%--------------------------------------------------------------------
\section{Internationalisation strategy}
\label{sec:i18n}

Internationalisation spans both the server and the client, so it is built
into the architecture rather than handled as a translation file added later.
English is the platform's primary language, with German maintained alongside
it from the start. The capability is present from the outset; only the
translations themselves are deferred, which keeps activating a new language
later cheap.

Human-readable text appears in 2 places, one per process. The single-page
application carries message catalogues, one per locale, for every label and
message; the backend renders its own outbound text, email above all, in
the recipient's locale. The two are connected by the locale itself: a user's
chosen locale is threaded through to the point where an email is rendered, so
the platform replies in the same language the user reads. Locale
resolution prefers an explicit per-user setting and falls back to English. The
chosen locale also governs date, number and currency formats, which differ
across the German-speaking markets; time-stamps are stored in UTC and rendered
in the user's zone.

%--------------------------------------------------------------------
\section*{Conclusion}
This chapter set out the architecture the rest of the report delivers. 5
drivers (auditability, integration-first design, EU data residency,
multi-tenant isolation and maintainability) led to CQRS and
Event~Sourcing, because together they make the compliance-grade audit trail
a property of the storage itself. The comby framework
supplies those patterns and the multi-tenant components, which is what lets
us cover the platform's requirements with only a small amount of custom code. The
requirements map onto a 4-context core domain model (Contract, Approval,
Deadline and DocumentLink, extended at its seams by the enterprise contexts)
on a single-backend topology hosted in the EU, with tenant isolation
and the audit trail as structural guarantees.

Chapter~\ref{chap:realisation} turns from this design to how it was
built, consolidating the engineering work delivered across the project:
the foundations, the contract core, and the enterprise capabilities.

%==============================================================================
\end{spacing}
